\documentclass{article}

\usepackage{amsmath,amssymb}
\usepackage{xurl}
\usepackage[hidelinks]{hyperref}

% Shared commands for notes in docs/paper-gaps/.

\DeclareUnicodeCharacter{2080}{\ifmmode{}_{0}\else\textsubscript{0}\fi}
\DeclareUnicodeCharacter{2081}{\ifmmode{}_{1}\else\textsubscript{1}\fi}
\DeclareUnicodeCharacter{2082}{\ifmmode{}_{2}\else\textsubscript{2}\fi}
\DeclareUnicodeCharacter{2099}{\ifmmode{}_{n}\else\textsubscript{n}\fi}

\newcommand{\Lean}{\textsc{Lean}}
\ExplSyntaxOn
\NewDocumentCommand{\leanid}{m}{
  \ifmmode
    \textsf{\detokenize{#1}}
  \else
    \group_begin:
    \urlstyle{sf}
    \tl_set:Nn \l_tmpa_tl {#1}
    \tl_replace_all:Nnn \l_tmpa_tl {\_} {_}
    \exp_args:NV \path \l_tmpa_tl
    \group_end:
  \fi
}
\ExplSyntaxOff
\providecommand{\tr}{\operatorname{tr}}
\newcommand{\tnleanIssueBase}{https://github.com/LionSR/TNLean/issues/}
\newcommand{\tnleanPrBase}{https://github.com/LionSR/TNLean/pull/}
\newcommand{\tnleanDocBase}{https://sirui-lu.com/TNLean/docs/}
\newcommand{\ghissue}[1]{\href{\tnleanIssueBase#1}{GitHub issue \##1}}
\newcommand{\ghpr}[1]{\href{\tnleanPrBase#1}{PR \##1}}
\newcommand{\doclink}[1]{\href{\tnleanDocBase#1.html}{\path{#1}}}

% Dirac notation and the identity operator, as in blueprint/src/macros/common.tex.
% \providecommand keeps note-local definitions authoritative.
\usepackage{dsfont}
\providecommand{\ket}[1]{|#1\rangle}
\providecommand{\bra}[1]{\langle #1|}
\providecommand{\braket}[2]{\langle #1 | #2 \rangle}
\providecommand{\ketbra}[2]{|#1\rangle\!\langle #2|}
\providecommand{\Id}{\mathds{1}}
\providecommand{\one}{\mathds{1}}
\providecommand{\C}{\mathbb{C}}
\providecommand{\MN}[1]{M_{#1}(\C)}
\providecommand{\MD}{M_D(\C)}

% External referencing for paper sources.  The build helper writes sanitized
% auxiliary files under build/paper-aux/ and excludes duplicated source labels.
\usepackage{xr}

% Import a generated paper auxiliary file with a caller-supplied label prefix.
% The candidate paths support builds from docs/paper-gaps/, the repository root,
% and build/paper-aux/ itself.
\newcommand{\importpaperaux}[2]{%
  \IfFileExists{../../build/paper-aux/#2.aux}{%
    \externaldocument[#1]{../../build/paper-aux/#2}%
  }{%
    \IfFileExists{build/paper-aux/#2.aux}{%
      \externaldocument[#1]{build/paper-aux/#2}%
    }{%
      \IfFileExists{#2.aux}{%
        \externaldocument[#1]{#2}%
      }{}%
    }%
  }%
}

% General external reference macros
\newcommand{\paperref}[2]{\ref{#1-#2}}
\newcommand{\papereqref}[2]{(\ref{#1-#2})}

% CPSV16 source (1606.00608 / MPDO-22-12-17-2.tex) external document and shortcuts
\importpaperaux{cpsv16-}{1606.00608/MPDO-22-12-17-2-tnlean}
\newcommand{\cpsvref}[1]{\paperref{cpsv16}{#1}}
\newcommand{\cpsveqref}[1]{\papereqref{cpsv16}{#1}}

% Verdict marker: \gapnote{<kind>}{<status>}. Kind is the policy
% classification (clarification, local-correction, scope-restriction,
% unfaithful, false-source, open-gap); status is open, wip (elimination
% underway), resolved, or historical. Machine-read by the site index;
% prints a verdict line.
\newcommand{\gapnote}[2]{\par\noindent\textbf{Verdict:} #1 (#2).\par}


\title{The Final Traced Subspin in the Definition of $\mathcal S_0$}
\author{}
\date{2026-07-12}

\begin{document}
\maketitle
\gapnote{local-correction}{resolved}

\section{The source sentence}

In the proof of Proposition~C.7 in Appendix~C.2 of
Ref.~\cite{Cirac2016MPDO_arXiv}, the map $\mathcal S_0$ reduces three physical
sites to their two outer subspins. The prose at line~1547 of
\path{Papers/1606.00608/MPDO-22-12-17-2.tex} states that $\mathcal S_0$ traces
the subspins $1r,2l,2r,2l$. The label $2l$ occurs twice, and the list contains
no subspin from the third site.

\section{Correction forced by the diagram}

For each site label $j$, let $L_j$ and $R_j$ denote its left and right subspin
spaces. In the diagram
\path{Papers/1606.00608/MPDO_TMatrix2.png} associated with
Ref.~\cite{Cirac2016MPDO_arXiv}, the three consecutive physical sector factors
before applying $\mathcal S_0$ are
\begin{align}
    (L_k\otimes R_k)
        \otimes(L_l\otimes R_l)
        \otimes(L_h\otimes R_h).
    \label{eq:cpgsv17_mpdo_s0_trace_subspin_typo_three_site_factors}
\end{align}
The site labels $k$, $l$, and $h$ correspond to the first, second, and third
sites, respectively. The map retains only the outer factors,
\begin{align}
    L_k\otimes R_h.
    \label{eq:cpgsv17_mpdo_s0_trace_subspin_typo_retained_outer_factors}
\end{align}
It therefore traces over
\begin{align}
    (R_k\otimes L_l)
        \otimes(R_l\otimes L_h).
    \label{eq:cpgsv17_mpdo_s0_trace_subspin_typo_traced_factors}
\end{align}
In site notation, the four factors in
Eq.~\ref{eq:cpgsv17_mpdo_s0_trace_subspin_typo_traced_factors} are the
subspins $1r,2l,2r,3l$. The displayed formula for $\mathcal S_2$ at
lines~1557--1559 of Ref.~\cite{Cirac2016MPDO_arXiv} retains the same two outer
factors. Hence the final occurrence of $2l$ at line~1547 is a typographical
error for $3l$.

For elements $l_j\in L_j$ and $r_j\in R_j$, the formal regrouping uses the
corrected equivalence
\begin{align}
    ((l_k,r_k),((l_l,r_l),(l_h,r_h)))
        &\longmapsto
        ((l_k,r_h),((r_k,l_l),(r_l,l_h))).
    \label{eq:cpgsv17_mpdo_s0_trace_subspin_typo_corrected_regrouping}
\end{align}
The corresponding formal declarations are recorded in
\doclink{TNLean/MPS/MPDO/RFPSubspinMaps}.\footnote{In particular,
\leanid{MPOTensor.ExplicitEtaOperators.s0RegroupEquiv},
\leanid{MPOTensor.ExplicitEtaOperators.s0Map}, and
\leanid{MPOTensor.ExplicitEtaOperators.s0SectorMap}.}
This correction changes no hypothesis and affects only the ordering of the
tensor factors.

\section{Verdict}

The repeated subspin label in the proof of Proposition~C.7 of
Ref.~\cite{Cirac2016MPDO_arXiv} requires a local correction. Replacing the
final $2l$ by $3l$ does not alter the proposition's statement or hypotheses.

\bibliographystyle{alpha}
\bibliography{references}

\end{document}
